
\section{Nuclear size: first-order energy shifts}
\textbf{Recognise:} a finite nucleus changes the electrostatic
potential only near \(r=0\), relative to the point-charge Coulomb model.

\textbf{Define unperturbed problem:}
\[
H_0\ket{nlm}=E_{nl}^{(0)}\ket{nlm},\qquad
\Delta V(r)=V_{\mathrm{finite}}(r)-V_{\mathrm{point}}(r).
\]
\textbf{First-order correction:}
\[
\Delta E_{nl}^{(1)}=\langle nlm|\Delta V|nlm\rangle
=\int_0^\infty |R_{nl}(r)|^2\Delta V(r)\,r^2\,\diff r.
\]
State the adopted nuclear charge model and radial normalisation.
If \(\Delta V\) is confined to a small nuclear radius,
inspect the small-\(r\) wavefunction: states with large density
at the nucleus feel the greatest correction.
\begin{methodbox}
\textbf{Method:} construct both potentials, subtract, set physical
integration limits, insert correctly normalised radial state,
integrate and sanity-check sign/units.
\end{methodbox}
